\begin{appendix}{Г}{информационное}
{Форматы сертификатов Государственной системы управления открытыми 
ключами}\label{NCES}

\section{Роли}\label{NCES.Roles}

В ГосСУОК прямо поддерживаются роли КУЦ, ИУЦ, ЦАС, ФЛ, ЮП и ЮЛ.
%
Роль ИУЦ проходит под именем <<РУЦ>>~--- республиканский удостоверяющий центр.
Это имя используется в настоящем приложении. Роль ЮП проходит как 
<<представитель юридического лица>>, роль ЮЛ~--- под именем <<ОРГ>>.

Сертификаты OCSP-сервера, СШВ, СЗД, TLS-сервера и КА дифференцируются только на 
уровне расширений, считаясь сертификатами единой роли TC (технологические 
сертификаты).
%
% skip: Наоборот, от ФЛ отделяется роль ФЛ2~--- физическое лицо как владелец 
% ID-карты. 

\section{Идентификационные данные}\label{NCES.Name}

Дополнительно к идентификационным атрибутам, указанным в
таблице~\ref{Table.ENTITIES.Attrs}, в сертификатах ГосСУОК используются
атрибуты, перечисленные в таблице~\ref{Table.NCES.Attrs}.

\begin{table}[H]
\caption{Дополнительные идентификационные атрибуты ГосСУОК}
\label{Table.NCES.Attrs}
\begin{tabular}{|l|p{7.2cm}|p{4.9cm}|}
\hline
Атрибут & Описание & Идентификатор\\
\hline
\hline
\code|description| & 
Текстовое описание объекта~\cite{RFC4519} &
\code|{2 5 4 13}|\\
\hline
%
\code|emailAddress| & 
Адрес электронной почты~\cite{RFC2985} &
\code|{1 2 840 113549 1 9 1}|\\
\hline
%
\texttt{идентификаторГИС} & 
Идентификатор государственной информационной системы &
Соответствует номеру, присвоенному в Государственном регистре
информационных систем\\
\hline                                      
\end{tabular}
\end{table}

В сертификате ЮП идентификаторы атрибутов \code{serialNumber} и 
\code{organizationalUnitName} меняются на \code|{1 2 112 1 2 1 1 1 1 1}| 
и \code|{1 2 112 1 2 1 1 5 2}| соответственно.

В таблице~\ref{Table.NCES.AttrRole} показано, какие атрибуты включаются в
идентификационные данные стороны ГосСУОК в зависимости от ее роли.
%
Знак <<$+$>> в таблице означает обязательное присутствие атрибута,
знак <<$\pm$>>~--- необязательное присутствие,
пропуск~--- обязательное отсутствие.

\begin{table}[hbt]
\caption{Идентификационные атрибуты ролей ГосСУОК}
\label{Table.NCES.AttrRole}
\begin{tabular}{|l|c|c|c|c|c|c|c|c|c|}
\hline 
Атрибут                       &КУЦ&РУЦ&ЦАС& ТС&  ФЛ & ЮП&ОРГ\\
\hline 
\hline 
\code{commonName}             & + & + & + & + &  +  & + & + \\
\code{surname}                &   &   &   &   &  +  & + &   \\
\code{name}                   &   &   &   &   &  +  & + &   \\
\code{givenName}              &   &   &   &   &  +  &   &   \\
\code{serialNumber}           &   &   &   &   &  +  & + &   \\
\code{countryName}            & + & + & + & + &  +  & + & + \\
\code{localityName}           &   & + & + & + &     & + & + \\
\code{stateOrProvinceName}    &   &   &   & + &     & + & + \\
\code{streetAddress}          &   &   &   & + &     & + & + \\
\code{organizationName}       & + & + & + & + &     &   & + \\
\code{organizationalUnitName} &   &   &   &   &     &   & + \\
\code{title}                  &   &   &   &   &     & + &   \\
\code{organizationIdentifier} &   & + & + &   &     &   & + \\
\hline                                                  
\code{description}            &   &   &   & + &     &   &   \\
\code{emailAddress}           &   &   &   & + &$\pm$& + &   \\
\texttt{идентификаторГИС}     &   &   &   & + &     &   &   \\
\hline
\end{tabular}
\end{table}

Имеются следующие особенности формирования идентификационных атрибутов.

\begin{enumerate}
\item
Ограничения на длины строковых атрибутов, указанные в 
таблице~\ref{Table.ENTITIES.Attrs}, ослабляются.

\item 
В сертификате ФЛ в атрибуте \code{commonName} указываются имя и 
фамилия физического лица на английском языке вместе с идентификационным  
(личным) номером из удостоверения лица.

\item 
В сертификате ФЛ в атрибутах \code{surname}, \code{name} и 
\code{givenName} указываются соответственно фамилия, имя и отчество 
физического лица на русском языке.

\item 
В сертификате ЮП в атрибуте \code{organizationName}
указывается полное название организации, а в атрибуте \code{name}~---
имя и отчество представителя организации на русском языке.
	
\item 
В сертификате ОРГ в атрибуте \code{organizationName}
указывается полное название организации, а в атрибуте \code{commonName}~---
сокращенное название.
\end{enumerate}

\section{Расширения сертификатов}\label{NCES.CertExt}

Дополнительно к расширениям сертификатов, описанным в~\ref{FMT.Ext}, в 
сертификатах ГосСУОК используются расширения, перечисленные в 
таблице~\ref{Table.NCES.CertExt}. 

\begin{table}[H]
\caption{Дополнительные расширения сертификатов ГосСУОК}
\label{Table.NCES.CertExt}
\begin{tabular}{|l|p{7.3cm}|p{5.3cm}|}
\hline
Расширение & Описание & Идентификатор\\
\hline
\hline
\code|SOAIdentifier| & 
Признак того, что субъект выступает в качестве источника 
привилегий (см. СТБ 34.101.67) & 
\code|{2 5 29 50}|\\
\hline
%
\texttt{УНП} & 
Учетный номер плательщика &
\code|{1 2 112 1 2 1 1 1 1 2}|\\
\hline
%
\texttt{УНПФ} & 
Учетный номер плательщика в органах Фонда социальной защиты населения & 
\code|{1 2 112 1 2 1 1 1 4 1}|\\
\hline                                      
\end{tabular}
\end{table}

В таблице~\ref{Table.NCES.Exts} показано, какие расширения включаются в
сертификат стороны ГосСУОК в зависимости от ее роли.
%
Знаки <<$+$>>, <<$\pm$>> и пропуски трактуются так же, как ранее.

\begin{table}[thb]
\caption{Расширения сертификатов ролей ГосСУОК}
\label{Table.NCES.Exts}
\begin{tabular}{|l|c|c|c|c|c|c|c|c|c|}
\hline 
Расширение                    &КУЦ&РУЦ&ЦАС& ТС& ФЛ  & ЮП  &  ОРГ\\
\hline                                    
\hline                                    
\code{SubjectAltName}         &   & + & + & + &  +  &     &  +  \\
\code{SubjectKeyIdentifier}   & + & + & + & + &  +  &  +  &  +  \\
\code{AuthorityKeyIdentifier} & + & + & + & + &  +  &  +  &  +  \\
\code{CertificatePolicies}    &   & + & + & + &$\pm$&  +  &  +  \\
\code{BasicConstraints}       & + & + & + & + &  +  &  +  &  +  \\
\code{CRLDistributionPoints}  &   & + & + & + &$\pm$&$\pm$&$\pm$\\
\code{AuthorityInfoAccess}    &   &   & + & + &  +  &  +  &  +  \\
\code{KeyUsage}               & + & + & + & + &  +  &  +  &  +  \\
\code{ExtKeyUsage}            &   &   &   & + &  +  &  +  &  +  \\
\hline                                    
\code{SOAIdentifier}          &   &   & + &   &     &     &     \\
\texttt{УНП}                  &   &   &   &   &     &  +  &  +  \\
\texttt{УНПФ}                 &   &   &   &   &     &     &  +  \\
\hline
\end{tabular}
\end{table}

Имеются следующие особенности формирования расширений.
\begin{enumerate}
\item 
В сертификате ФЛ в расширении \code{SubjectAltName}
указываются фамилия, имя и отчество субъекта на белорусском языке.
Фамилия задается в атрибуте \algname{surname-by} с идентификатором 
\code|{1 2 112 1 2 1 1 1 4 2}|, 
имя~--- в атрибуте \algname{givenName-by} с идентификатором 
\code|{1 2 112 1 2 1 1 1 4 3}|,
отчество~--- в атрибуте \algname{secondName-by} с идентификатором 
\code|{1 2 112 1 2 1 1 1 4 4}|. 
\item 
В расширении \code{CertificatePolicies} указываются идентификаторы политик 
ГосСУОК.                                                                   	
\item 
В сертификатах СШВ и СИ в расширении \code{KeyUsage} не указывается флаг 
\code{nonRepudiation}, а в сертификате КА, наоборот, указывается.
\item 
В сертификатах OCSP-сервера, СШВ и СЗД в расширении 
\code{KeyUsage} дополнительно указывается флаг \code{keyEncipherment}.
\item 
В сертификатах OCSP-сервера, СШВ, СЗД, СИ и TLS-сервера в расширении 
\code{KeyUsage} дополнительно указывается флаг \code{keyAgreement}.
\item 
В сертификате ЮП в расширении \code{KeyUsage} не указывается
флаг \code{nonRepudiation}, но дополнительно указываются флаги 
\code{keyAgreement} и \code{dataEncipherment}.
\end{enumerate}

\mbox{}

\end{appendix}
